\documentclass[11pt,addpoints]{exam}
\usepackage{amsfonts,amssymb,amsmath, amsthm, mathtools}
\usepackage{graphicx}
\usepackage{systeme}
\usepackage{pgf,tikz,pgfplots}
\pgfplotsset{compat=1.15}
\usepgfplotslibrary{fillbetween}
\usepackage{mathrsfs}
\usetikzlibrary{arrows}
\usetikzlibrary{calc}
\usepackage{moresize}
\usepackage{wasysym}
\usepackage{multicol}
\usepackage{enumitem}
\renewcommand{\thechoice}{{\Large$\Circle$ \ }}
\renewcommand{\choicelabel}{\thechoice}
\newcommand{\versionmark}{$\heartsuit$}
\renewcommand{\thepartno}{\roman{partno}}
\pagestyle{headandfoot}
\setlength\parindent{0pt}
\usepackage{array}
\usepackage{amssymb}
\checkboxchar{$\Box$}
\newcommand{\practiceShortSpace}{\par\vspace{0.55in}}
\newcommand{\practiceWorkSpace}{\par\vspace{0.85in}}

\firstpageheader{CS 2050 - Exam 4 Ryland Grace\\ Spring 2026}{}{Name: \underline{\hspace{2.5in}}\\ GTID: \underline{\hspace{2.5in}}}
\runningheader{CS 2050 - Exam 4 Ryland Grace}{Page \thepage\ of \numpages}{Name: \underline{\hspace{2 in}}}
\runningheadrule
\firstpagefooter{}{}{}
\runningfooter{}{}{}
\bracketedpoints
\pointsinmargin

\begin{document}

\begin{center}
\fbox{\fbox{\parbox{6in}{\centering
No notes, calculators, or other aids are allowed. Read all directions carefully and write your answers in the space provided.\\
Taking this exam signifies that you are aware of and acting in accordance with the Georgia Tech Academic Honor Code.\\ Do not separate any pages from the rest of your exam.
}}}
\end{center}
\vspace{\baselineskip}
\begin{center}
    \vspace{30mm}
    \HUGE{Version Ryland Grace}

    \numpoints\ points
\end{center}
\newpage

\textbf{For the following questions, completely fill in the bubble corresponding to your chosen answer.}

\begin{questions}

\question[2] Select the choice below to receive 2 free points.
\begin{choices}
    \choice I want free points!
\end{choices}
\practiceWorkSpace

\question[3] The Pigeonhole Principle states that if $n$ objects are placed into $m$ boxes and $m>n$, then at least one box must contain more than one object.
\begin{multicols}{2}
\begin{choices}
    \choice True
    \choice False
\end{choices}
\end{multicols}
\practiceShortSpace

\question[3] For integers $n\ge 1$ and $1\le k\le n$, the inequality
\[
\binom{n}{k}\le n^k
\]
holds.
\begin{multicols}{2}
\begin{choices}
    \choice True
    \choice False
\end{choices}
\end{multicols}
\practiceShortSpace

\question[3] There are more length-five binary strings that begin with $0$ than length-six binary strings that contain exactly three $0$s.
\begin{multicols}{2}
\begin{choices}
    \choice True
    \choice False
\end{choices}
\end{multicols}
\practiceShortSpace

\question[3] On the spacecraft \textit{Hail Mary}, a procedure consists of three successive tasks, $T_1,T_2,$ and $T_3$. For each $i$, task $T_i$ can be completed in exactly $n_i$ ways, regardless of the choices made for the other tasks. The procedure can therefore be completed in exactly $n_1n_2n_3$ ways.
\begin{multicols}{2}
\begin{choices}
    \choice True
    \choice False
\end{choices}
\end{multicols}
\practiceShortSpace

\question[3] For integers $n\ge 0$ and $0\le k\le n$, the number of $k$-element subsets of $\{1,2,\ldots,n\}$ equals the number of $(n-k)$-element subsets.
\begin{multicols}{2}
\begin{choices}
    \choice True
    \choice False
\end{choices}
\end{multicols}
\practiceShortSpace

\newpage

\question[5] Dr. Grace has 10 black, 10 red, and 10 blue markers. He takes 10 markers to class, including at least two of each color. What is the largest number of markers of one color that he is guaranteed to have?
\begin{choices}
    \choice 2
    \choice 4
    \choice 6
    \choice 10
    \choice None of the above.
\end{choices}
\practiceWorkSpace

\question[5] Grace assigns 25 indistinguishable tasks among 6 distinct Eridian engineers. In how many ways can he make the assignment if every engineer receives at least two tasks?
\begin{choices}
    \choice $\binom{25}{2}\binom{23}{2}\binom{21}{2}\cdots\binom{15}{2}$
    \choice $\binom{30}{5}$
    \choice $25^6$
    \choice $\binom{18}{5}$
    \choice None of the above.
\end{choices}
\practiceWorkSpace

\question[5] Dr. Ryland Grace's science exam has two versions, A and B. He must give one version to each of five distinct students receiving accommodations. If there are no restrictions on the assignments, in how many ways can he assign the versions?
\begin{choices}
    \choice 10
    \choice 25
    \choice 32
    \choice 60
    \choice None of the above.
\end{choices}
\practiceWorkSpace

\question[5] A four-member crew for the \textit{Hail Mary} is chosen from two scientists, three captains, and four engineers. If the crew must contain exactly one scientist and exactly one captain, how many crews are possible?
\begin{choices}
    \choice 24
    \choice 36
    \choice 48
    \choice 60
    \choice None of the above.
\end{choices}
\practiceWorkSpace

\newpage

\question[5] Aboard the \textit{Hail Mary}, Grace places 6 of 13 distinct books on the top level of a bookshelf and the remaining 7 books on the lower level. If the order of the books on each level matters, how many arrangements are possible?
\begin{choices}
    \choice $P(7,7)P(6,6)$
    \choice $\binom{13}{7}\cdot7\cdot6$
    \choice $P(13,7)P(6,6)$
    \choice $\binom{13}{7}\binom{13}{6}$
    \choice None of the above.
\end{choices}
\practiceWorkSpace

\question[5] How many distinct arrangements of the letters in \textsc{CIENCIA} have all vowels consecutive and all consonants consecutive? For example, \textsc{NCCIIEA} is valid, whereas \textsc{IECNCIA} is not.
\begin{choices}
    \choice $P(7,3)$
    \choice 36
    \choice $2P(7,3)P(7,4)$
    \choice 72
    \choice None of the above.
\end{choices}
\practiceWorkSpace

\question[5] Two distinct engineers and three distinct scientists sit at a circular table to discuss the \textit{Hail Mary}'s launch plans. How many seatings are possible if the engineers may not sit next to each other? Two seatings are considered the same exactly when one is a rotation of the other.
\begin{choices}
    \choice 6
    \choice 12
    \choice 24
    \choice 120
    \choice None of the above.
\end{choices}
\practiceWorkSpace

\question[5] Grace has 15 identical ``I Love Earth'' pins to distribute among 4 distinct crewmates, one of whom is the captain. If every crewmate must receive at least 2 pins and the captain may receive at most 5 pins, how many distributions are possible?
\begin{choices}
    \choice $\binom{10}{3}$
    \choice $\binom{10}{3}-\binom{6}{3}$
    \choice $\binom{18}{3}-\binom{12}{3}$
    \choice $\binom{10}{3}-\binom{5}{3}$
    \choice None of the above.
\end{choices}
\practiceWorkSpace

\newpage
For the following question, show \textbf{all} work using topics covered in lecture. Brute-force solutions receive no credit. Simplify the final result to a single integer and write it on the answer line. Leaving the answer in combinatorial notation may result in a point penalty.

\question[12] Consider ordered sequences of length five whose entries belong to $\{-1,0,1\}$. How many such sequences have entries that sum to zero? \textbf{Show all work, simplify your answer to a single integer, and write it on the answer line below.}

\vspace{\stretch{1}}
Answer: \rule{4cm}{0.15mm} \hspace{1cm}

\newpage
\question[16] Aboard a United Nations ship at sea, 170 scientists are working on \textit{Project Hail Mary}. Prove that at least three scientists have birthdays that fall in the same month of the year and on the same day of the week (for example, both birthdays occur on a Monday in October).

\vspace{\stretch{1}}
\newpage
\question[15] Let $n$ and $k$ be integers with $2\le k\le n$. Give a double-counting proof of the identity
\[
\binom{n}{2}=\binom{k}{2}+k(n-k)+\binom{n-k}{2}.
\]

\vspace{\stretch{1}}
\newpage
\question \textbf{EXTRA CREDIT:} Art contest! Draw a picture of a CS 2050 TA and write the TA's name. Submissions will be judged on a scale from 1 to 3. You may not raise your hand to ask a TA anything about this question.

\vspace{\stretch{1}}
\newpage
\textbf{Scratch Paper:} If you use this page for a solution, clearly indicate both here and on the original question page that the solution appears here.

\end{questions}
\end{document}
